% !TeX root = PROJECT_ALIGNMENT_BOARD.tex
% !TeX program = pdflatex
\documentclass[10pt,a4paper]{article}
\usepackage[T1]{fontenc}
\usepackage[utf8]{inputenc}
\usepackage{lmodern}
\usepackage[margin=20mm]{geometry}
\usepackage{titlesec}
\usepackage{enumitem}
\usepackage{tabularx}
\usepackage{booktabs}
\usepackage{amssymb}
\usepackage{microtype}
\usepackage[hidelinks]{hyperref}
\setlength{\parindent}{0pt}
\setlength{\parskip}{2pt}
\titleformat{\section}{\normalsize\bfseries}{}{0pt}{}
\titlespacing*{\section}{0pt}{8pt}{3pt}
\setlist[itemize]{leftmargin=*,itemsep=1pt,topsep=2pt,parsep=0pt}
\newcommand{\blank}{\rule{3cm}{0.4pt}}
\newcommand{\checkitem}{\item[$\square$]}
\begin{document}

\begin{center}
{\LARGE\bfseries Group Q Project Alignment Board}\par
\smallskip
{\large Research Roots}\par
\smallskip
14 October 2026
\end{center}

\textbf{Purpose:} Agree on one product for the report, presentation and video.
The direction below is recorded; four-member confirmation remains pending.

\textbf{Members:} Yehao Chen, Peng Wang, Dennis de Hoon and Gerald Tobenna Eluagu.

\section*{1. Product and Scope}
\begin{itemize}
\item A website for students and researchers comparing planned research with existing work.
\item Core output: one conceptual knowledge-landscape circle and two trees.
\item Similarity supports borrowing methods, datasets, protocols and research resources.
\item Distinctiveness identifies differences and possible directions for further investigation.
\item Tags help others discover reusable work, subject to permissions and sharing choices.
\item Earlier self-assessment calculations, three-module navigation and a third tree are outside the current core scope.
\end{itemize}

\section*{2. Result Meaning}
\begin{itemize}
\item \textbf{Circle:} A conceptual map of the research landscape; its labels and placement need agreement. It does not measure the boundary of all human knowledge.
\item \textbf{Similarity Tree:} Related topics, questions, methods and resources, with sources and reasons for each connection.
\item \textbf{Distinctiveness Tree:} Differences from retrieved work, with evidence and possible research questions. Differences do not prove originality.
\item \textbf{Frontier directions:} Questions within the distinctiveness view, rather than a separate tree or an unvalidated frontier score.
\item \textbf{Sustainability:} Potential benefits through reuse and efficient resource use; measured savings have not been demonstrated.
\end{itemize}

\section*{3. Pages and User Flow}
\begin{itemize}
\item \textbf{Input page:} Research question, method and relevant data/equipment needs; review the description before searching.
\item \textbf{Results page:} Circle, Similarity Tree and Distinctiveness Tree; expand branches for source records, explanations and reuse opportunities.
\item \textbf{User control:} Revise input, inspect sources and flag misleading connections.
\item \textbf{Optional sharing:} Review tags and visibility before sharing. Decide whether this belongs in the demo or future work.
\end{itemize}

\section*{4. Sources and Responsibilities}
\begin{itemize}
\item Candidate sources: TU/e research-portal records and OpenAlex; record actual access, coverage and permissions.
\item AI may assist retrieval, comparison, tagging and explanation; record tools and prompts actually used.
\item Researchers judge relevance, reuse feasibility and research direction.
\item Show missing information and coverage limits; check explanations against original sources.
\item Ethical priorities: explicability, privacy, fairness and researcher control.
\end{itemize}

\newpage
\section*{5. Decisions to Confirm}
\begin{itemize}
\checkitem One agreed product statement and target user.
\checkitem Meaning and labels of the circle and both trees.
\checkitem Qualitative descriptions or validated numerical indicators.
\checkitem One shared sample input, source set and output for the demo.
\checkitem Tags, sharing options and features included in this version.
\checkitem Owners of pages, report sections, slides and recording.
\end{itemize}

\textbf{Agreed product statement:} \hrulefill\par
\textbf{Chosen demonstration:} \hrulefill\par
\textbf{Open disagreements:} \hrulefill

\section*{6. Team Contributions}
Assignments below are blank until accepted by the members.

\begin{tabularx}{\textwidth}{@{}lXX@{}}
\toprule
\textbf{Member} & \textbf{Accepted task/output} & \textbf{Target date/status} \\
\midrule
Yehao Chen & & \\[8pt]
Peng Wang & & \\[8pt]
Dennis de Hoon & & \\[8pt]
Gerald Tobenna Eluagu & & \\[8pt]
\bottomrule
\end{tabularx}

\section*{7. Report Presentation and Video}
\begin{itemize}
\item Use the same product, page labels, example and claims in all deliverables.
\item Align B and F with research matching; update C from three trees to two; make G show the agreed interface.
\item Complete H with actual prompts, sources and tools; keep unfinished work visibly incomplete.
\item Pitch sequence: problem and user; website; shared demonstration; reuse benefits; limitations and next evaluation step.
\item Label illustrative outputs and distinguish proposed functions from tested results.
\item Separate pitch: 3--5 minutes according to the Phase 2 format. Verify the class-presentation length and deadlines on Canvas.
\end{itemize}

\section*{8. Discussion and Approval}
For each new item, record:
\begin{itemize}
\item Topic and proposal.
\item Member comments and supporting source/file.
\item Agreed decision, reason, owner and target date.
\item Status: Open / Agreed / In progress / Checked.
\end{itemize}

\textbf{Meeting date:} \blank\quad\textbf{Shared document link:} \blank

\textbf{Member confirmation:}\par
Yehao \blank\quad Peng \blank\par
Dennis \blank\quad Gerald \blank

\smallskip
\textit{Use one shared editable copy. Record disagreements before finalising the report or recording.}

\end{document}
